\subsection{Ressourcen- und Zeitplanung}\label{sec:ressourcen}

Die Projektarbeit umfasste 5~ECTS, entsprechend 150~Stunden, verteilt auf vier Wochen konzeptionelle Arbeit (160~Stunden Brutto-Kapazität bei einer wöchentlichen Kapazität von 40~Stunden) und eine fünfte Woche für die Berichterstellung. Die Gesamtkapazität inklusive Bericht belief sich auf rund 200~Stunden.

\begin{figure}[ht!]
    \centering
    \includegraphics[width=0.95\textwidth]{images/2.1_Ressourcen-und-Zeitplanung/Zeitplan.png}
    \caption[Ressourcen- und Zeitplanung]{Ressourcen- und Zeitplanung des UI-Redesign-Projekts. Quelle: Eigene Darstellung.}
    \label{fig:zeitplan}
\end{figure}

Die Ressourcenallokation priorisierte die Analyse- und Konzeptionsphase. Jede Woche beinhaltete feste Zeitblöcke für interne Synchronisation. Eine Reflexion der initialen Planung im Vergleich zur tatsächlichen Ressourcenallokation erfolgt in \autoref{sec:reflexion_fazit}. Als gestalterische UI-Ressourcen wurden das Vaadin-Aura-Theme, Lucide Icons und Instrument Sans eingesetzt. Die konkrete Verwendung ist im StyleTile (\autoref{sec:styletile}) dokumentiert. Die im Projekt eingesetzten Werkzeuge sind in \hyperref[sec:anhang_werkzeuge]{Anhang C} dokumentiert.
